\documentclass[11pt]{article}
\usepackage[a4paper,margin=28mm]{geometry}
\usepackage[T1]{fontenc}
\usepackage[utf8]{inputenc}
\usepackage{lmodern}
\usepackage{amsmath,amssymb,amsthm,mathtools}
\usepackage{booktabs,array}
\usepackage{xcolor}
\usepackage{microtype}
\usepackage[hidelinks]{hyperref}
\setlength{\emergencystretch}{3em}

\definecolor{deepblue}{RGB}{24,65,108}
\hypersetup{
  colorlinks=true,
  linkcolor=deepblue,
  citecolor=deepblue,
  urlcolor=deepblue,
  pdftitle={The Reconstructed Theory Has One Mass: A Machine-Checked Volume-Uniform Transfer Gap for a Finite Ising Strip},
  pdfauthor={Lluis Eriksson},
  pdfsubject={Replacement of ai.viXra.org:2607.0078 (v1)}
}
\newtheorem{theorem}{Theorem}[section]
\newtheorem{proposition}[theorem]{Proposition}
\newtheorem{remark}[theorem]{Remark}
\newcommand{\RR}{\mathbb{R}}
\newcommand{\CC}{\mathbb{C}}
\newcommand{\cH}{\mathcal{H}}
\newcommand{\vac}{\Omega}
\newcommand{\ip}[2]{\langle #1,#2\rangle}
\newcommand{\norm}[1]{\lVert #1\rVert}

\title{\textbf{The Reconstructed Theory Has One Mass}\\[4pt]
\large A machine-checked volume-uniform transfer gap for a finite Ising strip}
\author{Lluis Eriksson}
\date{4 August 2026\\[4pt]
\small Replacement of \href{https://www.ai.vixra.org/abs/2607.0078}{ai.viXra.org:2607.0078} (v1)}

\begin{document}
\maketitle

\begin{abstract}
We give a finite-dimensional Osterwalder--Schrader (OS) reconstruction for an
anisotropic Ising model on open rectangular strips and prove a spectral gap
whose rate is uniform in the spatial extent.  For temporal coupling $\beta$,
spatial coupling $\gamma$, and
\[
  0<\alpha<1,\qquad
  2\tanh |\beta|+2\tanh |\gamma|\leq \alpha,
\]
one number $m>0$ works for every spatial length.  Reflection positivity is
proved directly for the Gibbs measure; the null space of the reflected form
is identified with the kernel of an explicit boundary-collapse map; the OS
quotient is linearly equivalent to the finite boundary-vector space; and the
transfer operator forced by the site and bond forms is intertwined with a
symmetrised transfer matrix.  Dobrushin comparison supplies positive Perron
data and the bound
\[
 \norm{T_L-|\vac_L\rangle\langle\vac_L|}\leq e^{-m}
 \quad\text{for every }L.
\]
Consequently, connected reconstructed correlations decay at the same common
rate.  Lean 4 checks the composition without assuming reflection positivity,
Perron data, a vacuum, a gap, or clustering at the public endpoint.  A
conservative priority survey found earlier machine-checked OS axioms and
abstract OS reconstruction, and earlier volume-uniform measure-side
clustering, but no earlier artifact combining a concrete reconstructed
transfer operator with a proved volume-uniform spectral gap.
\end{abstract}

\section{Why the composition matters}

Reflection positivity and high-temperature mixing normally live in different
mathematical descriptions.  The first is a positivity statement about a
reflected Euclidean measure and produces a physical Hilbert space and a
transfer operator.  The second is a contraction statement about conditional
distributions and produces decay of measure-level correlations.  A paper that
merely places those facts side by side has not shown that the reconstructed
operator inherits the quantitative rate.  Conversely, an abstract structure
containing a field named ``gap'' does not prove a gap for a model.

The result here is the missing composition in a finite system.  All spaces are
finite, but the spatial dimension of the slice grows with $L$, and the
conclusion uses one $m$ before quantifying over $L$.  Thus the statement is not
the tautology that every fixed positive matrix has some finite-volume gap.
It is also not a thermodynamic-limit construction or a continuum Yang--Mills
mass-gap theorem; the precise boundary of the claim is discussed in
Section~\ref{sec:limits}.

This manuscript supersedes---rather than supplements---the three internal
stepping-stone manuscripts on the reconstructed gap, correlator decay, and
OS-forced transfer operator.

This is the promised \emph{commutant plus bulk dissipation} candidate in its
classical reduction.  On the reconstructed boundary space, multiplication by
functions on $\Sigma_L$ is a maximal abelian algebra, hence is its own
commutant.  The vacuum line is the invariant scalar sector, while the
orthogonal complement is the bulk sector.  The estimate proved below is
exactly uniform dissipation of that complement.  Nothing here promotes this
classical statement to a noncommutative commutant theorem; doing so would be a
separate paper and is deliberately not claimed.

The formal development is in
\path{YangMills/OS/OSReconstructionUniform.lean}. It composes earlier,
independently checked modules rather than replacing them with a monolithic
proof.  This makes the logical inputs visible and permits a direct
non-vacuity audit.

\section{The finite anisotropic Ising strip}

Let
\[
  \Sigma_L=\{0,1\}^{\{0,\ldots,L\}}
\]
be the spatial slice space, with the two spin values interpreted as $\pm1$.
The open-chain spatial weight is
\[
  w_{\gamma,L}(\sigma)
   =\prod_{j=0}^{L-1}\exp\!\bigl(\gamma s(\sigma_j)s(\sigma_{j+1})\bigr)>0,
\]
and the bond between consecutive time slices is
\[
  K_{\beta,L}(\sigma,\tau)
   =\prod_{j=0}^{L}\exp\!\bigl(\beta s(\sigma_j)s(\tau_j)\bigr).
\]
For a path $X=(X_0,\ldots,X_N)$ of spatial slices, the unnormalised Gibbs
weight is
\[
 W_{L,N}(X)=\prod_{t=0}^{N}w_{\gamma,L}(X_t)
             \prod_{t=0}^{N-1}K_{\beta,L}(X_t,X_{t+1}).
\]
These definitions contain no operator.  The model-to-operator dictionary is
therefore a theorem rather than a definitional identification.

The symmetrised kernel is
\[
 S_L(\sigma,\tau)=
 \sqrt{w_{\gamma,L}(\sigma)}\,K_{\beta,L}(\sigma,\tau)
 \sqrt{w_{\gamma,L}(\tau)}.
\]
If $\lambda_L>0$ is its Perron value, define the normalised kernel
$M_L=S_L/\lambda_L$.  Its Euclidean continuous linear operator is denoted
$T_L$.  For a strictly positive Perron vector $\Omega_L$, its Euclidean
normalisation is denoted $\widehat\Omega_L$.

\section{The three load-bearing inputs}

\subsection{Reflection positivity on the measure}

For half-path observables $F_i$ and coefficients $c_i\in\CC$, the formal
theorem \path{os_ising_measure_gram_nonneg} proves that the full Gibbs
sum
\[
 \sum_{i,j}\overline{c_i}c_j
 \sum_X \overline{F_i(X_{\rm past})}F_j(X_{\rm future}^{\rm rev})W_{L,N}(X)
\]
is a nonnegative real number when $\beta\geq0$.  This is a statement about the
measure itself, not only about an auxiliary quadratic form.

\subsection{The OS quotient is identified}

The boundary-collapse map sums a half-path observable against its Gibbs
weight while retaining the terminal slice.  The theorem
\path{os_ising_null_space_iff} states that its kernel is exactly the
null space of the reflected site form.  Surjectivity then gives the bundled
linear equivalence
\[
  \frac{\{\text{half-path observables}\}}
       {\ker(\text{collapse})}
  \simeq_{\CC} \CC^{\Sigma_L},
\]
implemented as \path{osIsingPhysicalEquiv}. Hence no unidentified quotient
is hidden in the reconstruction.

The transfer operator $\mathcal T_L$ on boundary vectors is forced by
\[
 \ip{u}{\mathcal T_Lv}_{\rm site}=\ip{u}{v}_{\rm bond}.
\]
The theorem \path{os_ising_transfer_is_measure_sum} identifies this
matrix element with the reflected Gibbs two-point sum one time step farther
apart.  Thus the operator conclusion has a measure-level reading.

\subsection{A uniform Dobrushin feed}

The rectangular interaction has row sum bounded by
$2\tanh|\beta|+2\tanh|\gamma|$.  The previously proved theorem
\path{dobrushin_ising_uniform_gap} converts the common Dobrushin window
into one $m>0$ and, for every $L$, positive Perron data satisfying
\[
  M_L\Omega_L=\Omega_L,
  \qquad
  \norm{T_L-|\widehat\Omega_L\rangle
                    \langle\widehat\Omega_L|}\le e^{-m}.
\]
The order of quantifiers is essential: $m$ is selected before $L$.

\section{Main theorem}

\begin{theorem}[Machine-checked OS reconstruction with one mass]
\label{thm:main}
Let $0<\alpha<1$, $\beta\geq0$, and
\[
  2\tanh|\beta|+2\tanh|\gamma|\leq\alpha.
\]
Then the finite Gibbs measure is reflection positive at every reflection
depth.  Its OS null space is the kernel of the explicit collapse map and its
physical quotient is linearly equivalent to $\CC^{\Sigma_L}$.  Moreover,
there exists $m>0$ such that for every $L$ there are
$\lambda_L>0$ and $\Omega_L>0$ for which:
\begin{enumerate}
\item $M_L\Omega_L=\Omega_L$;
\item $(T_L,\widehat\Omega_L)$ is a symmetric unit-vacuum transfer pair;
\item
$\norm{T_L-|\widehat\Omega_L\rangle\langle\widehat\Omega_L|}
 \le e^{-m}$;
\item for every real boundary vector $v$ and $n\in\mathbb N$,
\[
 \left|\ip{v}{T_L^n v}
   -\ip{\widehat\Omega_L}{v}^2\right|
 \le \norm{v}^2 e^{-mn};
\]
\item multiplication by $\sqrt{w_{\gamma,L}}$ intertwines the transfer
operator forced by the OS forms with the unnormalised symmetrised kernel
$S_L$.
\end{enumerate}
\end{theorem}

In Lean, the quantitative portion is the single declaration
\path{os_reconstruction_uniform_gap}. Its hypotheses are exactly the
three scalar inequalities $0<\alpha$, $\alpha<1$, and the Dobrushin window.
The condition $\beta\geq0$ belongs only to the separate reflection-positivity
theorem.  In particular, the endpoint does not take any of items (1)--(4) as
an argument.

\begin{remark}[Meaning of ``one mass'']
This means one $m>0$ chosen before $L$ and valid for every finite $L$; it
asserts neither a unique particle excitation nor an infinite-volume
Hamiltonian or isolated mass shell.
\end{remark}

\section{Proof architecture}

The formal proof is short because the interfaces are strong; its content is
the exact compatibility of those interfaces.

First, the Dobrushin corollary produces $m$, and for each $L$ produces
$\lambda_L$, $\Omega_L$, positivity, the fixed-point equation, and the
projected norm bound.  Symmetry of $M_L$, strict positivity of $\Omega_L$,
and the fixed-point equation are passed to
\path{vacuumTransfer_opOf}. That theorem proves, entrywise, symmetry of
the operator, unit norm of $\widehat\Omega_L$, and
$T_L\widehat\Omega_L=\widehat\Omega_L$.

Second, \path{clustering_of_gap} applies the projected norm estimate to
the identity
\[
 \ip{v}{T_L^n v}-\ip{\widehat\Omega_L}{v}^2
   =\ip{v}{(T_L-P_{\widehat\Omega_L})^n v}
 \quad(n\ge1),
\]
with the $n=0$ case treated by Cauchy--Schwarz.  This yields the same $m$,
not a second rate selected after $L$.

Third, the existing similarity theorem
\path{transferOp_sqrtWeightEquiv} gives
\[
 \mathcal T_L Q_L=Q_L S_L,
 \qquad (Q_Lu)(\sigma)=\sqrt{w_{\gamma,L}(\sigma)}u(\sigma).
\]
Strict positivity of $w_{\gamma,L}$ makes $Q_L$ invertible.  The operator
whose projected norm is controlled is therefore the normalised version of
the same symmetrised transfer matrix that is intertwined with the OS-forced
operator.

Finally, the two measure-identification theorems attach that operator back to
the reflected Gibbs sums.  This last step is what rules out an isolated
finite-matrix result with no reconstructed interpretation.

\section{Adversarial non-vacuity audit}

Two failure modes were tested directly.

\paragraph{Empty antecedent.}
The window is inhabited: for example $\beta=\gamma=0.1$ and $\alpha=0.5$
satisfy it.  More generally, continuity of $\tanh$ gives a nonempty open
high-temperature neighbourhood of the origin.  Every Boltzmann weight is a
strictly positive exponential, so partition functions and the normalising
Perron mass are nonzero.

\paragraph{Conclusion smuggled into the hypotheses.}
The public endpoint contains no argument of type \texttt{VacuumTransfer}, no
Perron vector, no spectral-norm inequality, and no clustering bound.  Its
only hypotheses are scalar parameter inequalities.  Reflection positivity is
also not postulated: the Gram theorem consumes only $\beta\ge0$ and constructs
the nonnegative real witness for the actual Gibbs sum.

The quotient is not circular either.  Its null subspace is proved equal to
the kernel of a concrete surjective linear map before the quotient
equivalence is formed.  Likewise, the transfer operator is not selected to
have a convenient spectrum; it is uniquely forced by the site/bond equation
and only afterwards intertwined with $S_L$.

\section{Why the general Birkhoff route was rejected}

A tempting route is to seek volume-uniform contraction in the standard
positive cone using the projective diameter of
$S_L=\sqrt w\,K\sqrt w$.  A pre-registered numerical falsifier evaluated
all pairwise Hilbert-metric ratios for $L=2,\ldots,8$, both inside and outside
the Dobrushin window.  It passed 88 of 88 internal consistency checks in both
ordinary and optimised Python modes and found
\[
 \Delta(S_L)=\Delta(K_L)=4|\beta|(L+1),
 \qquad
 \kappa_L=\tanh\!\bigl(|\beta|(L+1)\bigr)\longrightarrow1.
\]
The equality of projective diameters is structural: positive diagonal
congruence does not change the relevant column cross-ratios.  Hence the slice
weight cannot rescue standard-cone Birkhoff contraction, and tensor
additivity makes the obstruction grow with $L$.  This does not exclude a
genuinely different cone, but it kills the proposed general Birkhoff proof.
The paper therefore uses the surviving Dobrushin reduction and makes no
beyond-window claim.

\section{Formal artefacts and reuse map}

Table~\ref{tab:files} records the load-bearing in-repository chain.  All paths
are relative to the repository root.

\begin{table}[ht]
\centering
\small
\begin{tabular}{>{\raggedright\arraybackslash}p{0.31\linewidth}
                >{\raggedright\arraybackslash}p{0.26\linewidth}
                >{\raggedright\arraybackslash}p{0.33\linewidth}}
\toprule
File & Reused declarations & Role here\\
\midrule
\path{YangMills/OS/SpatialOS.lean} &
\path{gibbsSum_reflected_gram_nonneg} & Gibbs-measure RP\\
\path{YangMills/OS/SpatialReconstruction.lean} &
\path{mem_ker_collapseL_iff}, \path{physicalEquiv},
\path{osPairing_transfer_gibbsSum}, similarity theorem &
quotient, forced operator, measure dictionary\\
\path{YangMills/OS/DobrushinCorollary.lean} &
\path{dobrushin_ising_uniform_gap}, \path{sliceW} &
one mass and model family\\
\path{YangMills/OS/DobrushinTransport.lean} &
\path{vacuumTransfer_opOf} & vacuum packaging\\
\path{YangMills/OS/TransferGap.lean} &
\path{clustering_of_gap} & operator gap to decay\\
\path{YangMills/OS/OSReconstructionUniform.lean} &
five public declarations & new composition endpoint\\
\bottomrule
\end{tabular}
\caption{Load-bearing formal reuse.}
\label{tab:files}
\end{table}

Five satellite repositories were inspected at their public heads on
4 August 2026.  Their role is negative or terminological, not load-bearing:
\begin{center}
\small
\begin{tabular}{lll}
\toprule
Repository & inspected HEAD & reusable verdict\\
\midrule
\path{lean-os-positivity} & \path{2ea7692f} & RP vocabulary; mother tree stronger\\
\path{lean-transfer-matrix} & \path{4e2366a4} & 1D Ising gap only\\
\path{lean-gaussian-field} & \path{53af447b} & finite covariance algebra only\\
\path{aqft-split-inclusion-series} & \path{078c39ae} & TeX/numerics, no Lean\\
\path{physmath-lean-lemmas} & \path{b26ce4c4} & unrelated summability closure\\
\bottomrule
\end{tabular}
\end{center}
No theorem was vendored from these repositories.  This is deliberate: source
presence is not treated as a materialised object file or as verification.

\section{Related machine-checked work and priority boundary}

Douglas, Hoback, Mei, and Nissim formalised the Glimm--Jaffe/OS axioms for the
four-dimensional Gaussian free field in Lean~4~\cite{douglas2026}.  This
already rules out any claim of first machine-checked OS axioms or reflection
positivity. The repository \path{mrdouglasny/reflection-positivity}
constructs an abstract physical Hilbert space and a self-adjoint contractive
transfer operator from reflection-positive data~\cite{reflectionrepo}; its
\path{GappedTransfer} interface carries the gap as a hypothesis rather than
proving it for a concrete family. The repository \path{mrdouglasny/lgt}
proves Dobrushin-style, volume-uniform measure-side correlation decay for
lattice gauge theory~\cite{lgtrepo}, but has no OS reconstruction or transfer
operator.  The active $P(\phi)_2$ project contains transfer-matrix spectral
infrastructure, with its current headline uniform estimates still represented
by named project assumptions~\cite{pphi2repo}.

Accordingly, the narrow priority statement is: the public survey found no
earlier machine-checked result combining (i) a concrete interacting Gibbs
family, (ii) an OS-reconstructed transfer operator, (iii) a gap proved rather
than assumed, and (iv) one rate uniform in spatial volume.  This is an
evidence-qualified statement, not a claim that abstract OS reconstruction,
reflection positivity, Dobrushin uniqueness, or finite-matrix spectral gaps
are new.

The classical foundations are the OS reconstruction papers
\cite{os1973,os1975}, the lattice reflection-positive construction of
Osterwalder and Seiler~\cite{osterwalderseiler1978}, and Dobrushin's
comparison method~\cite{dobrushin1968}.  The present work is a finite,
machine-checked quantitative composition of these themes.

\section{Reproducibility and audit status}

The repository is pinned to Lean \path{v4.29.0-rc6} and Mathlib commit
\begin{center}
\path{07642720480157414db592fa85b626dafb71355b}.
\end{center}
The OS-R numerical gate contains 110 registered checks and is run with Python assertions enabled and
disabled.  Lean compilation, full core compilation, and the repository-wide
axiom oracle are executed only on a Colab Pro+ CPU high-RAM runtime.  The
pre-measurement predictions are 8481 core jobs and 3015 oracle reports.
On the exact frozen merged tree the realised totals are 8482 core jobs and
3015 oracle reports, with zero \texttt{sorryAx}; the incorrect core prediction
is retained in the record rather than repaired after measurement.  The final
fresh-clone run ended with every child sentinel and the atomic parent sentinel
equal to zero.

The final frozen commit, byte counts, LF and CRLF hashes, realised build
counts, and Colab opening/closing times are recorded in the accompanying
verification ledger and release manifest.  This paper does not self-certify
the terminal verdict: it supplies the frozen object and the attacks attempted;
an independent audit task must issue the final assessment.

\section{Limitations and next targets}
\label{sec:limits}

The theorem concerns finite open rectangles and uniformity in the spatial
extent.  It does not construct an infinite-volume measure, a continuum limit,
a relativistic Wightman theory, or a Hamiltonian by functional calculus.  The
correlation theorem is stated for real Euclidean boundary vectors.  Complex
observables can be split into real and imaginary parts, but that corollary is
not needed for the headline endpoint.

Uniform finite-volume bounds alone do not produce an infinite-volume transfer
operator: one still needs compatible maps between boundary spaces, convergence
of states and operators, and independence of finite-volume boundary
conditions.  These inputs are not present here.

The Dobrushin window is sufficient and deliberately non-sharp.  The failed
standard-cone Birkhoff test explains why a naive tensor-product improvement
cannot enlarge it uniformly.  A credible beyond-window project would need a
new cone, a block dynamics with a proved uniform contraction, or a comparison
principle that survives the tensorial wall.  None is claimed here.

\section*{Conclusion}

The main formal fact is compact: one scalar high-temperature window produces
one mass before the spatial volume is chosen, and the resulting operator is
the one reconstructed from the reflected Gibbs measure.  Keeping positivity,
quotient identification, similarity, Perron normalisation, and decay in the
same theorem suite closes the two most dangerous loopholes: a vacuous
implication and a gap smuggled into the assumptions.  The result is modest in
physical scope but, within that scope, fully quantitative and mechanically
checkable.

\begin{thebibliography}{10}

\bibitem{os1973}
K.~Osterwalder and R.~Schrader.
\newblock Axioms for Euclidean Green's functions.
\newblock \emph{Communications in Mathematical Physics}, 31:83--112, 1973.
\newblock \url{https://projecteuclid.org/euclid.cmp/1103858969}.

\bibitem{os1975}
K.~Osterwalder and R.~Schrader.
\newblock Axioms for Euclidean Green's functions II.
\newblock \emph{Communications in Mathematical Physics}, 42:281--305, 1975.
\newblock \url{https://projecteuclid.org/euclid.cmp/1103899050}.

\bibitem{osterwalderseiler1978}
K.~Osterwalder and E.~Seiler.
\newblock Gauge field theories on a lattice.
\newblock \emph{Annals of Physics}, 110:440--471, 1978.
\newblock DOI: \href{https://doi.org/10.1016/0003-4916(78)90039-8}
{10.1016/0003-4916(78)90039-8}.

\bibitem{dobrushin1968}
R.~L. Dobrushin.
\newblock The description of a random field by means of conditional
probabilities and conditions of its regularity.
\newblock \emph{Theory of Probability and Its Applications}, 13:197--224,
1968.
\newblock DOI: \href{https://doi.org/10.1137/1113026}{10.1137/1113026}.

\bibitem{douglas2026}
M.~R. Douglas, S.~Hoback, A.~Mei, and R.~Nissim.
\newblock Formalization of QFT.
\newblock arXiv:2603.15770, 2026.
\newblock \url{https://arxiv.org/abs/2603.15770}.

\bibitem{reflectionrepo}
M.~R. Douglas and contributors.
\newblock \texttt{reflection-positivity}: reflection positivity and transfer
operators in Lean 4, 2026.
\newblock \url{https://github.com/mrdouglasny/reflection-positivity}.

\bibitem{lgtrepo}
M.~R. Douglas and contributors.
\newblock \texttt{lgt}: lattice gauge theory and Dobrushin comparison in
Lean 4, 2026.
\newblock \url{https://github.com/mrdouglasny/lgt}.

\bibitem{pphi2repo}
M.~R. Douglas and contributors.
\newblock \texttt{pphi2}: construction of $P(\phi)_2$ quantum field theory
in Lean 4, 2026.
\newblock \url{https://github.com/mrdouglasny/pphi2}.

\end{thebibliography}

\end{document}
